\newpage
\begin{frame}
	\renewcommand{\figurename}{Figure}
	\fta{Nodal analysis}

	\s{
		Analysing an electrical network can sometimes be a complex task. As the network grows, 
        so does the complexity of the analysis. In some cases, the analysis methods presented so far for node 
        and mesh analysis are not sufficient to determine all the variables of an electrical network. In this case, the 
        node potential method offers a possible solution for analysis. In the nodal analysis, all other potentials are determined based on a 
        reference potential, with an assignment of 0V. Using the node potential method
        , the electrical network from Figure \ref{BildKnotenpotentialverfahren} is analysed as follows. 
    }

	\b{
		\begin{itemize}
			\item Preparation of the network
			\item Determination of nodes and node potentials
			\item Allocation of source currents
			\item Setting up the conduct value matrix
			\item Setting up the system of equations
		\end{itemize}
	}

	\fu{   
        \resizebox{0.7\textwidth}{!}{%                                         
        \input{Tikz/src/knotenpotentialverfahren.tex}
        }                                                  
    }{{\bf Network for Nodal analysis.} Electrical network with a voltage source, a current source and four resistors. The network is used to 
    explain the Nodal analysis. \label{BildKnotenpotentialverfahren}}  
	
\end{frame}


\newvideofile{Knotpot}{Knotenpotentialverfahren}
\begin{frame}
	\ftx{Learning objectives Node potential method}
	\begin{Learning objectives}{Node analysis}
        The Students
        \begin{itemize}
			\item are familiar with the steps of the nodal analysis
            \item can analyse electrical networks using the nodal analysis method.
        \end{itemize}                   
    \end{Learning objectives}

	\speech{KnotpotLernz}{1}{Das Knotenpotentialverfahren ist neben dem Maschenstromverfahren eine Möglichkeit, um elektrische Netzwerke zu analysieren.
	Hier sollen die notwengien Schritte zur Anwendung des Knotenpotentialverfahrens vorgestellt werden.} 
\end{frame}

\begin{frame}
	\renewcommand{\figurename}{Figure}
	\ftb{Preparation of the network}

	\s{
		The nodal analysis method works with current sources and conductance values. In addition to the current source, the network shown has 
        a voltage source and resistance values. These need to be converted. As shown in Figure \ref{BildKnot1}, the voltage source is transformed into a current source. The series resistance of the 
        voltage source becomes a parallel resistance for the current source.
	}

	\onslide<1->{
	\b{
		\begin{itemize}
			\item Conversion of the voltage source into a current source:
		\end{itemize}
	}

	\fu{  
        \resizebox{0.7\textwidth}{!}{%                                     
         \input{Tikz/src/knot1.tex}
		}                                               
    }{{\bf Conversion of a voltage source.} Conversion of a voltage source with series resistance to a current source with 
    parallel resistance. \label{BildKnot1}}}   

	\s{
		When converting voltage sources into current sources, the voltage values of the voltage sources must also be converted into 
		current values. This is done according to equation \ref{GleichungUmrechnungSpannungsquelle}. 
	}

	\onslide<2->{
	\b{
		\begin{itemize}
			\item Conversion of voltage source to current source:
		\end{itemize}
	}

	\begin{eq}
			I_\mathrm{q} = U_\mathrm{q} \cdot G_\mathrm{i} \label{GleichungUmrechnungSpannungsquelle}
		\end{eq}}

	\s{
		Since the nodal analysis method does not use the resistance values of the components, but rather the conductance values,
        all resistance values are converted into conductance values according to equation \ref{GleichungUmrechnungWiderstand}.
	}	

	\onslide<3->{
	\b{
		\begin{itemize}
			\item Calculation of the conductance values:
		\end{itemize}
	}

    \begin{eq}               
         G_\mathrm{i} = \frac{1}{R_\mathrm{i}} \label{GleichungUmrechnungWiderstand}
    \end{eq}}

	\s{
		Figure \ref{BildKnot2} shows the conversion of the voltage sources from the specified network into equivalent 
		current sources. The voltage source $U_0$ and the resistor $R_0$ become the current source $I_\mathrm{q0}$ and 
 		the parallel conductance $G_0$. The remaining resistors are also converted into conductances.
	
		\fu{   
        \resizebox{\textwidth}{!}{%                                         
			\input{Tikz/src/knot2.tex}
        }                                                  
    }{{\bf Conversion of the network for the nodal analysis method.} The voltage source $U_0$ becomes the current source $I_\mathrm{q0}$, and the 
    resistance values R become conductances G. \label{BildKnot2}}  
	}

	\speech{KnotpotVorb}{1}{Beim Knotenpotentialverfahren wird mit Stromquellen gearbeitet. 
	Deshalb wird nun gezeigt wie Spannungsquellen in die benötigten Stromquellen umgewandelt werden können. 
	Es bietet sich an, zuvor noch das elektrische Netzwerk zu vereinfachen, beispielsweise durch die Zusammenfassung von Parallelschaltungen.
	Eine Spannungsquelle mit Innenwiderstand in Serie wird wie gezeigt in eine Stromquelle mit parallelem Innenwiderstand umgewandelt.} 
	\speech{KnotpotVorb}{2}{Für die Umwandlung muss der Spannungswert der Spannungsquelle noch in den Stromwert für die Stromquelle überführt werden.
	Dies erfolgt über das Ohmsche Gesetz.}
	\speech{KnotpotVorb}{3}{Die Leitwerte für das Knotenpotentialverfahren ergeben sich aus dem Kehrwert der Widerstandswerte.} 
\end{frame}


\begin{frame}
    \ftx{Nodal analysis}

    \begin{Key point}{Nodal analysis}
        In the nodal analysis method, an electrical network with current sources and conductances is analysed. 
    \end{Key point}

	\speech{KnotpotMrk}{1}{Beim Knotenpotentialverfahren wird ein elektrisches Netzwerk mit Stromquellen und Leitwerten analysiert.} 
\end{frame}


\begin{frame}
	\renewcommand{\figurename}{Figure}
	\ftb{Determination of nodes and node potentials}

	\s{
		Now the nodes must be determined. To do this, we define a reference node with the index 0, i.e. $K_0$.
		Next, all nodes are numbered consecutively. The electrical network with the numbered nodes is shown in
		Figure \ref{BildKnot3}. The nodes drawn without components between the numbered nodes
        have the same potential and thus form a node.  
	}

	\onslide<1->{
	\b{
		\begin{itemize}
			\item Determination of the reference node $K_0$ and the consecutive nodes $K_1$ and $K_2$:
		\end{itemize}
	}
		
	\fu{   
		\resizebox{0.7\textwidth}{!}{%                                         
		\input{Tikz/src/knot3.tex}
		}                                                  
	}{{\bf Converted network.} Determination of the reference node $K_0$ and the consecutive nodes $K_1$ and $K_2$. \label{BildKnot3}}} 

	\s{
		To determine the node potentials, all potentials of nodes $K_1$ and $K_2$ are referenced to the reference node.
        $k-1$ equations are required to perform the node potential method. Referring to the specified network, 
		this results in two equations, which are written in vector notation according to equation \ref{GleichungKnotenpotentiale}.
		Here, the potential for node $K_1$ is defined by the voltage from node $K_1$ to the reference node.	
		This results in the notation $U_\mathrm{K1-K0}$. The same is done for node $K_2$. Here, the voltage 
        $U_\mathrm{K2-K0}$ is defined.
	}

	\onslide<2->{
	\b{
		\begin{itemize}
			\item $k-1$ equations are required. This results in the following stress vector:
		\end{itemize}
	}
		
	\begin{eq}
        U_\mathrm{K} = 
        \begin{bmatrix}
            U_\mathrm{K1-K0} \\
            \\
            U_\mathrm{K2-K0} \\
        \end{bmatrix} \label{GleichungKnotenpotentiale}
    \end{eq}}

	\speech{KnotpotKnot}{1}{Als weiteren Schritt müssen die Knoten gekennzeichnet werden und ein Bezugsknoten festgelegt werden.
	Der Bezugsknoten wird normalerweise mit null gekennzeichnet und die folgenden Knoten fortlaufend durchnummeriert.
	Hier können somit drei Knoten definiert werden.} 
	\speech{KnotpotKnot}{2}{Für das Knotenpotentialverfahren müssen k minus eins Gleichungen aufgestellt werden. 
	So werden hier zwei Spannungsgleichungen definiert, welche in vektorieller schreibweise notiert werden.
	Die Spannungen geben den potentialunterschied vom jeweiligen Knoten zum Bezugsknoten an.} 
\end{frame}

\begin{frame}
	\ftb{Allocation of source currents}

	\s{
		For the specified nodes, apart from the reference node, the outgoing and incoming currents must be recorded.
		Each current source at the nodes is noted. Here, incoming current sources are given a positive sign and outgoing
		current sources are given a negative sign. 
		The current from the current source $-I_\mathrm{q0}$ flows out of node $K_1$, and this 
		current is assigned a negative sign. The current from the current source $-I_\mathrm{q2}$ flows out of node $K_2$, and 
		this is also assigned a negative sign. According to equation \ref{GleichungQuellströme}, this results in the 
        vector for the source currents.
 
	}

	\b{
		\begin{itemize}
			\item Assignment of source currents for each node
			\item Outflowing spring streams have a negative sign.
			\item Inflowing source streams are noted positively.
		\end{itemize}
	}

	\begin{eq}
        I_\mathrm{K} = 
        \begin{bmatrix}
            -I_\mathrm{q0} \\
            \\
            -I_\mathrm{q2} \\
        \end{bmatrix} \label{GleichungQuellströme}
    \end{eq}

	\speech{KnotpotQuell}{1}{Folgend werden die Quellströme notiert. Dabei werden abgesehen vom Bezugsknoten für jeden Knoten
	die Stromquellen festgehalten, die in den Knoten einfließen oder ausfließen. Beim Knoten Ka eins fließt der Strom der Quelle 
	I Q null raus. Aus dem Knoten Ka zwei fließt der Strom der Quelle I Q zwei raus. Da es sich bei beiden Knoten und beiden Strömen
	um abfließende Ströme handelt, werden sie mit einem negativen Vorzeichen versehen.} 
\end{frame}

\begin{frame}
	\renewcommand{\figurename}{Figure}
	\ftb{Conductance matrix}

	\s{
        After converting the resistance values into conductance values, the conductance matrix is created using the conductance values. To do this, 
        the nodes apart from the reference node are examined again. For this purpose, the electrical network presented is shown once more in 
        Figure \ref{BildLeitwertmatrix} with coloured branches.
 		The conductance values of the nodes are noted along the main diagonal. 
		For each node, the conductance values of the directly adjacent components are written down. 
        For node $K_1$, these would be the conductance values $G_0$, $G_1$ and $G_3$ adjacent to the branch highlighted in blue, and 
        for node $K_2$, the conductance values $G_2$ and $G_3$. Away from the main diagonal, the conductance values that lie directly 
		between the nodes are recorded on the other elements of the conductance matrix. 
		Between nodes $K_1$ and $K_2$, the branch highlighted in green with the conductance value $G_3$ lies directly between them. 
		If 	there are other nodes between the nodes under consideration, 
        these are not directly connected to each other and a 0 is entered in the element of the conductance matrix. The conductance matrix created in this 
        way is illustrated in equation \ref{GleichungLeitwertmatrix}. 
	}

	\b{
		\begin{itemize}
			\item<2-> Main diagonal with the conductance values directly adjacent to the nodes
			\item<3-> Other elements with the conductance values that lie directly between the nodes
		\end{itemize}
	}

	\fu{
		\resizebox{0.6\textwidth}{!}{%  
		\input{Tikz/src/leitwertmatrix.tex}		
		}
	}{{\bf Electrical network with highlighted connections.} Creation of the conductance matrix for the nodal analysis method 
	\label{BildLeitwertmatrix}}

	\b{
		\only<1>{
		\begin{eq}
			G_\mathrm{K} =
			\begin{bmatrix} 
				0 & 0 \\
				\\
				0 & 0 \\
			\end{bmatrix} 
		\end{eq}
		}
		\only<2>{
		\begin{eq}
			G_\mathrm{K} =
			\begin{bmatrix} 
				\color{blue}{G_0 + G_1 + G_3} & 0 \\
				\\
				0 & G_2 + G_3 \\
			\end{bmatrix} 
		\end{eq}
		}
		\only<3>{
		\begin{eq}
			G_\mathrm{K} =
			\begin{bmatrix} 
				\color{blue}{G_0 + G_1 + G_3} & \color{green}{-G_3} \\
				\\
				\color{green}{-G_3} & G_2 + G_3 \\
			\end{bmatrix} 
		\end{eq}
		}
	}

	\s{
		\begin{eq}
			G_\mathrm{K} =
			\begin{bmatrix} 
				\color{blue}{G_0 + G_1 + G_3} & \color{green}{-G_3} \\
				\\
				\color{green}{-G_3} & G_2 + G_3 \\
			\end{bmatrix} 
			\label{GleichungLeitwertmatrix}
		\end{eq}
	}

	\speech{KnotpotLeit}{1}{Als nächstes wird die Leitwertmatrix aufgestellt. Es handelt sich bei diesem Netzwerk um eine zwei mal zwei Matrix.} 
	\speech{KnotpotLeit}{2}{Auf der Hauptdiagonalen der Matrix werden die Leitwerte eingetragen, welche direkt an den jeweiligen Knoten angrenzen.
	Für den Knoten Ka eins sind das beispielsweise die Leitwerte G null, G eins und G drei.} 
	\speech{KnotpotLeit}{3}{Die Übrigen Elemente der Matrix werden mit den LEitwerten gefüllt, welche direkt zwischen den Knoten liegen.
	Da sich in diesem Beispiel nur der Leitwert G drei direkt zwischen den knoten befindet, wird dieser zwischen den Knoten eingetragen.
	Die Leitwerte, welche sich zwischen den Knoten befinden, werden in der Matrix mit einem negativen Vorzeichen versehen.} 
\end{frame}

\begin{frame}
	\ftb{Set up a system of equations}

	\s{
		After the vector of node potentials, the vector of source currents, and the conductance matrix have been determined, these equations are used to 
        set up the system of equations for the node potential method, which is shown in Equation \ref{GleichungKnot1}.
}

	\b{
		\begin{itemize}
			\item Solving the system of equations using Gauss's method or Cramer's rule
		\end{itemize}
	}

	\begin{eq}
        \begin{bmatrix} 
            G_0 + G_1 + G_3 & -G_3 \\
            \\
            -G_3 & G_2 + G_3 \\
        \end{bmatrix}
		\cdot
		\begin{bmatrix}
            U_\mathrm{K1-K0} \\
            \\
            U_\mathrm{K2-K0} \\
        \end{bmatrix}
		=
		\begin{bmatrix}
            -I_\mathrm{q0} \\
            \\
            -I_\mathrm{q2} \\
        \end{bmatrix}
    \end{eq} \label{GleichungKnot1}

	\s{        
		The system of equations of the node potential method can then be solved using the Gauss elimination method 
        as an example. By rearranging the equation according to the node potentials sought, these can be identified. If, for example,
 		the voltage across resistor $R_1$ has been calculated in the network under investigation, the current through 
        the resistor can be determined. As an alternative to the Gaussian elimination method, Cramér's rule can also be used
        to solve linear systems of equations. 
	}

	\speech{KnotpotGleich}{1}{Mithilfe der Leitwertmatrix, dem Spannungsvektor und dem Vektor der Quellströme lässt sich nun 
	das angegebene Gleichungssystem aufstellen. Durch das Umstellen der Gleichung nach den gesuchten Knotenpotentialen können 
	diese identifiziert werden. Um lineare Gleichungsysteme zu lösen, kann beispielsweise das Gaußsche Eliminationsverfahren oder 
	die Kramersche Regel verwendet werden.} 
\newpage

\end{frame}

\begin{frame}
	\renewcommand{\figurename}{Figure}

    \ftx{Example: Nodal analysis}
    \begin{expl}{Node potential method}{}

    \s{
        Analysis of the electrical network from Figure \ref{BildKnot4} using the node potential method. 
    }

	\begin{enumerate}
        \only<1>{			
            \item[]
    The following steps should be carried out for the node potential method:

    \begin{itemize}
        \item[a)] Preparation of the network
		\item[b)] Determination of nodes and node potentials
		\item[c)] Allocation of source currents
		\item[d)] Preparation of the conductivity matrix
		\item[e)] Set up a system of equations
    \end{itemize}

    \fu{  
        \resizebox{0.5\textwidth}{!}{%                                     
        \input{Tikz/src/knot4.tex}
        }                  
    }{{\bf Example.} Node potential method on an electrical network. \label{BildKnot4}}                                            
}


        \only<2>{			
            \item[a)] Preparing the network:
			\begin{figure}[H]
                \resizebox{0.65\textwidth}{!}{%
                    \input{Tikz/src/knot5.tex}	
                }\centering                                                           
            \end{figure}
            }
        \only<3>{
            \item[b)] Determination of nodes and node potentials:
            \onslide<3>{
			\begin{figure}[H]
             	\resizebox{0.4\textwidth}{!}{%                                                        
                    \input{Tikz/src/knot6.tex}
                }\centering                                                           
            \end{figure}
            \pause

			\begin{eq}
				U_\mathrm{K} = 
				\begin{bmatrix}
					U_\mathrm{K1-K0} \\
					\\
					U_\mathrm{K2-K0} \\
					\\
					U_\mathrm{K3-K0} \\
				\end{bmatrix} \nonumber
			\end{eq}
			}
        }
        \only<4>{
            \item[c)] Allocation of source currents:
            \onslide<4>{
				\begin{eq}
					I_\mathrm{K} = 
					\begin{bmatrix}
						0 \\
						\\
						-I_\mathrm{q0}+I_\mathrm{q1} \\
						\\
						-I_\mathrm{q1}-I_\mathrm{q2} \\
					\end{bmatrix} \nonumber
				\end{eq}
            }
        }
        \only<5>{
            \item[d)] Determination of the conductance matrix:
            \onslide<5>{
                \begin{eq}
					G_\mathrm{K} =
					\begin{bmatrix} 
						G_3 + G_4 + G_5 & -G_3 & -G_5 \\
						\\
						-G_3 & G_0 + G_1 + G_3 & -G_1 \\
						\\
						-G_5 & -G_1 & G_1 + G_2 + G_5 \\
					\end{bmatrix} \nonumber
				\end{eq}
            }
        }
        \only<6>{
            \item[e)] Set up a system of equations:
            \onslide<6>{
                \begin{eq}
                    \begin{bmatrix} 
                        G_3 + G_4 + G_5 & -G_3 & -G_5 \\
						\\
						-G_3 & G_0 + G_1 + G_3 & -G_1 \\
						\\
						-G_5 & -G_1 & G_1 + G_2 + G_5 \\
                    \end{bmatrix}   
                    \cdot
                    \begin{bmatrix}
                        U_\mathrm{K1-K0} \\
						\\
						U_\mathrm{K2-K0} \\
						\\
						U_\mathrm{K3-K0} \\
                    \end{bmatrix}
                    =
                    \begin{bmatrix}
                        0 \\
						\\
						-I_\mathrm{q0}+I_\mathrm{q1} \\
						\\
						-I_\mathrm{q1}-I_\mathrm{q2} \\
                    \end{bmatrix}\nonumber
                \end{eq} 
            }
        }
    \end{enumerate}    
\end{expl}{}{}

\speech{KnotpotBsp}{1}{Die Schritte zur Anwendung des Knotenpotentialverfahrens wurden nun erläutert. 
Folgend wird anhand des abgebildeten Beispielnetzwerkes noch einmal die grundsätzliche Arbeitsweise wiederholt.
Hierzu wird das Netzwerk enstprechend für das Knotenpotentialverfahren vorbereitet, die Knoten und die Knotenpotentiale bestimmt,
die Quellströme zugeordnet, die Leitwertmatrix und anschließend das Gleichungssystem aufgestellt.} 
\speech{KnotpotBsp}{2}{Da das Beispielnetzwerk nicht ausschließlich aus Stromquellen besteht, 
müssen die Spannungsquellen mit Serienwiderständen in Stromquellen mit Parallelwiderständen umgewandelt werden.} 
\speech{KnotpotBsp}{3}{Anhand des umgewandelten Netzwerkes werden die Knoten und die zugehörigen Potential definiert.
Hieraus ergibt sich U K der Vektor der Knotenpotentiale.} 
\speech{KnotpotBsp}{4}{Der Stromvektor der Quellströme ergibt sich entsprechend zu I K.
Hier erfolgen keine Einströmungen oder Ausströmungen für den Knoten K Eins. Es wird eine 0 notiert.
Beim Knoten K 2 fließt der Strom I Q 0 aus dem Knoten heraus und der Strom I Q 1 in den Knoten K 2 hinein.
Aus dem Knoten K 3 fließen die Ströme I Q 1 und I Q 2 heraus.} 
\speech{KnotpotBsp}{5}{Die Leitwertmatrix des Beispielnetzwerkes wird als G K dargestellt.} 
\speech{KnotpotBsp}{6}{Zum Schluss wird das Gleichungssystem aufgestellt. Nach den gesuchten Größen umgestellt und berechnet.} 
\end{frame}